\section{Conclusion}
\label{sec:ch5-conclusion}

This paper explored how biophilic experiences can inform the design of future interactive indoor environments beyond prevailing comfort and optimisation models. Through three design sessions with architects and designers, we contribute what is, to our knowledge, the first biophilic interaction design framework for HCI and HBI: a mapped design space that renders experiential configurations comparable, design values that surface the normative orientations shaping biophilic translation, and biophilic interaction forms that reframe how interaction with indoor environments can occur over time. By bridging architectural imaginaries with interaction design research, this work provides a foundation for cumulative, comparative, and generative inquiry into technologically mediated buildings that engage ecological processes, multisensory experience, and forms of interaction that de-centre direct human control and align with ecological orientations.